% ===================================================================
%  કોષ્ટક નં. ૧૫ — બજાર. --- ખાવાની વસ્તુઓ.
%
%  Traduction, faite en 2026, en gujarati d'aujourd'hui, sur l'ido de
%  texto/io. Regles de la colonne : voir 10-kostak-01.tex. Une seule
%  numerotation. 28 blocs, 102 renvois — le plus dense du livret.
%
%  QUATRIEME FOIS QUE CETTE COLONNE ECRIT CE QU'ON NE MANGE PAS ICI,
%  ET CETTE FOIS CE N'EST PLUS UN INTERDIT RELIGIEUX, C'EST UN CODE
%  PENAL. Le cochon du tableau 7 etait un interdit de caste et de
%  religion ; le menu du tableau 13 etait de la coutume ; la viande de
%  bœuf (94) est autre chose. L'abattage de la vache est puni au
%  Gujarat de la reclusion a perpetuite depuis 2017 — c'est la peine
%  la plus lourde de toute l'Inde. Le mot existe et sert tous les
%  jours : « ગૌમાંસ » est dans la loi, dans les jugements, dans les
%  journaux. La boucherie de la planche en vend ; la colonne l'ecrit,
%  sans note et sans detour, comme elle a ecrit ભૂંડ et le
%  Saint-Emilion. LA LANGUE A LES MOTS ; C'EST LE FAIRE QUI EST
%  INTERDIT, PAS LE NOMMER — et il fallait qu'un tableau le verifie
%  la ou l'interdit n'est pas seulement dans les tetes.
%
%  LE MEME ALINEA CONTINUE : l'ail (101) et l'oignon, que la planche
%  fait vendre par un gamin, sont refuses absolument par les jains et
%  par une grande partie des vaishnavas gujaratis — non parce qu'ils
%  sont des betes, mais parce qu'on les tient pour echauffants. Une
%  cuisine entiere, celle des jours de jeune, se fait sans eux. Les
%  mots sont pourtant les plus ordinaires de la langue : લસણ, ડુંગળી.
%
%  L'EPICERIE (32) EST L'ENDROIT DU LIVRET OU LES DEUX BOUTS D'UNE
%  MEME ROUTE SE TOUCHENT. Le mot francais dit l'epice, et l'epice
%  est venue d'ici : Cambay, Bharuch et Surat ont charge les navires
%  qui ont fini par faire du mot epicier un nom de metier en Europe.
%  Le mot gujarati, « કરિયાણાની દુકાન », dit la meme chose par l'autre
%  bout — કરિયાણું, ce sont les denrees seches, celles qui se
%  transportent. Et dans la vitrine de cette epicerie-la, l'ido
%  enumere le the, le chocolat, le sucre, l'ananas et les bananes :
%  tout ce qui a voyage.
%
%  ET TROIS MOTS PORTUGAIS DE PLUS, TOUS LES TROIS DANS L'ASSIETTE.
%  Le releve en comptait sept au tableau 11 — મેજ la table, પાદરી le
%  pretre, અનનાસ l'ananas, મિસ્ત્રી le macon, બંબો la pompe, ઇસ્પિતાલ
%  l'hopital — et un huitieme au tableau 14, આયા la bonne d'enfants.
%  Ce tableau-ci en apporte trois d'un coup : « બટાટા » la pomme de
%  terre (61), de batata ; « કોબી » le chou (17), de couve ;
%  « પાંઉ » le pain (79), de pão. DIX EN QUINZE TABLEAUX. Et les trois
%  derniers ne sont pas des mots d'administration ou de chantier : ce
%  sont la pomme de terre et le chou, sans lesquels une cuisine
%  gujaratie d'aujourd'hui ne tient pas debout, et le pain du
%  vada-pav, qui est le plat de rue le plus mange de l'ouest de
%  l'Inde. LES PORTUGAIS N'ONT PAS SEULEMENT APPORTE DES MOTS : ILS
%  ONT APPORTE LES LEGUMES, ET LE MOT EST VENU ACCROCHE A LA CHOSE.
%
%  Deux mots pour finir, qui n'avaient pas encore servi. « હરાજી »
%  (24) est la criee, un mot de negoce que le Gujarat emploie du
%  marche aux legumes a la salle des ventes. Et « ભરતકામ » (85) est la
%  broderie : celle du Kutch et du Saurashtra, au miroir et au fil de
%  soie, qu'on connait dans le monde entier sous ce nom-la. La
%  brodeuse de la planche travaille pour une lingere de
%  sous-prefecture ; le mot, lui, est celui d'un metier qui n'a jamais
%  cesse.
% ===================================================================

%%K t15-apar-1 apar
\VUsaut{4.0mm}
\VUtitre{15.0pt}{60}{કોષ્ટક નં. ૧૫}
\VUsaut{3.0mm}

%%K t15-tit-1 sub
\VUpk{11.4pt}{40}{બજાર. --- ખાવાની વસ્તુઓ.}
\VUsaut{3.0mm}

%%K t15-01-1 p
1. --- અમારા ખોરાક માટે જરૂરી માલ અમને પહોંચાડનારા મોટા ભાગના વેપારીઓ
\VUgras{ઢાંકેલા બજારના} \VUgras{મંડપ} \textsuperscript{(1)} નીચે કે
પડોશના રસ્તાઓમાં ભેગા થાય છે.

%%K t15-02-1 p
2. --- \VUgras{ભૂંડના માંસનો વેપારી} \textsuperscript{(2)}, જે
\VUgras{હૅમ} \textsuperscript{(3)}, \VUgras{લોહીની સૉસેજ}
\textsuperscript{(4)}, \VUgras{સૉસેજ} \textsuperscript{(5)},
\VUgras{આંતરડાંની સૉસેજ} \textsuperscript{(6)} અને \VUgras{ચરબી}
\textsuperscript{(7)} વેચે છે, તે \VUgras{માખણ-ચીઝ વેચનારી બહેન}
\textsuperscript{(8)} પાસે ઊભો રહે છે, જેની માંડણી લાલ પોપડાવાળાં
\VUgras{ડચ ચીઝ} \textsuperscript{(9)}, \VUgras{કૅમેમ્બેર ચીઝ}
\textsuperscript{(10)}, \VUgras{ગ્રુયેરનાં ચક્કર}
\textsuperscript{(11)} અને તાજા \VUgras{માખણના લોંદાથી}
\textsuperscript{(12)} ભરેલી છે.

%%K t15-03-1 p
3. --- તેની આગળ એક \VUgras{શાકવાળી બહેન} \textsuperscript{(13)}
પોતાની ગ્રાહક બહેનોને સુંદર \VUgras{ટમેટાં} \textsuperscript{(14)},
\VUgras{ફ્લાવર} \textsuperscript{(15)}, \VUgras{સલાડનાં પાન}
\textsuperscript{(16)}, \VUgras{કોબી} \textsuperscript{(17)},
\VUgras{કોળાં} \textsuperscript{(19)}, \VUgras{ટેટી}
\textsuperscript{(18)} અને \VUgras{બિલાડીના ટોપ} \textsuperscript{(20)}
સુધ્ધાં આપે છે. પણ એક \VUgras{બિયર પહોંચાડનારે}, જેણે પોતાની
\VUgras{માલ પહોંચાડવાની ગાડી} \textsuperscript{(21)} બરાબર તેની માંડણી
આગળ ઊભી રાખી છે --- જે \VUgras{બિયરનાં પીપ} \textsuperscript{(22)} થી
ભરેલી છે ---, તેના ગ્રાહકોને પાસે આવતાં રોક્યા છે. એક \VUgras{માળણ} તેને
\VUgras{આર્ટિચોકનો} \textsuperscript{(23)} ટોપલો લાવી આપે છે.

%%K t15-tit-2 sub
\VUsaut{4.0mm}
\VUpk{10.2pt}{40}{વેચવું અને ખરીદવું.}
\VUsaut{3.0mm}

%%K t15-04-1 p
4. --- બજારમાં ભારે ચહલપહલ છે, કેમ કે \VUgras{હરાજીનો}
\textsuperscript{(24)} સમય થયો છે. \VUgras{ગૃહિણીઓ}
\textsuperscript{(26)}, જે પોતાની ખરીદી માટે ત્યાં આવે છે, તે
\VUgras{આંતરડાં વેચનારી બહેનની} \textsuperscript{(25)} દુકાન આગળથી
પસાર થતાં થોભે છે, \VUgras{હોજરીનાં આંતરડાં} કે \VUgras{વાછરડાનું
માથું} ખરીદવા.

%%K t15-05-1 p
5. --- એક જાડો \VUgras{રસોઇયો} \textsuperscript{(27)} \VUgras{માછલી
વેચનારી બહેન} \textsuperscript{(28)} પાસે ખૂબ સુંદર \VUgras{ટુનાનો}
ભાવ કરે છે, જે પોતાની મરજી પ્રમાણે ભાવ ચડાવ્યે જાય છે. તેને
\VUgras{ઈલ, સોલ માછલી, ટ્રાઉટ} અને \VUgras{કાર્પનો} મોટો જથ્થો મળ્યો
છે. તેની \VUgras{કૉડ} અને \VUgras{છીપલાં} ખૂબ તાજાં છે.

%%K t15-tit-3 sub
\VUsaut{4.0mm}
\VUpk{10.2pt}{40}{સસ્તો અને મોંઘો માલ.}
\VUsaut{3.0mm}

%%K t15-06-1 p
6. --- \VUgras{મરઘાં વેચનારી બહેન} \textsuperscript{(29)}, જે તેની
પાસે બેઠી છે, તે ખૂબ પ્રામાણિક છે. જ્યારે તે \VUgras{ટર્કી, હંસ,
બતક} કે સાદું \VUgras{મરઘીનું બચ્ચું} વેચે છે, ત્યારે તે વાજબી ભાવ જ
માગે છે.

%%K t15-07-1 p
7. --- ચોકના ખૂણે, પહેલા માળે, થોડા વખતથી એક \VUgras{કાપડ વેચનારો}
\textsuperscript{(30)} બેઠો છે, જેની પાસે ખરીદનારી બહેનોને
\VUgras{મેજપોશ, નૅપકિન, એપ્રન} અને \VUgras{લૂછવાનાં કપડાંનો} ખૂબ
સુંદર જથ્થો મળવાની હંમેશાં ખાતરી હોય છે. એ જ માળે એક
\VUgras{છરીવાળાએ} \textsuperscript{(31)} પોતાની દુકાન બેસાડી છે. તે
મંડપની વેચનારીઓની \VUgras{છરીઓ} ધારદાર કરી આપે છે અને માળીઓ માટે સૌથી
કઠણ \VUgras{પોલાદનાં} \VUgras{છોલવાનાં દાતરડાં} બનાવે છે.

%%K t15-tit-4 sub
\VUsaut{4.0mm}
\VUpk{10.2pt}{40}{કરિયાણું.}
\VUsaut{3.0mm}

%%K t15-08-1 p
8. --- તેની દુકાનની નીચે થોડા વખતથી ખાવાની વસ્તુઓની એક મોટી દુકાન
ખૂલી છે. એ હવે પહેલાંની સાદી અને નજીવી \VUgras{કરિયાણાની દુકાન}
\textsuperscript{(32)} રહી નથી, જે માત્ર રોજિંદો માલ વેચતી —
\VUgras{ચા} \textsuperscript{(33)}, \VUgras{ચૉકલેટ}
\textsuperscript{(34)}, \VUgras{ગાંગડા ખાંડ} \textsuperscript{(37)}
અને \VUgras{સાબુ} \textsuperscript{(38)}. ત્યાં હવે ગમે તે વસ્તુ મળે
છે — \VUgras{કડક બિસ્કિટ}, \VUgras{ડબાબંધ ખોરાક}
\textsuperscript{(35)}, \VUgras{લિકર} \textsuperscript{(36)},
\VUgras{રમ, કૉન્યાક, કિર્શ, વરિયાળીનો દારૂ} અને
\VUgras{ક્યુરાસાઓ}. ત્યાં શિકારનું માંસ પણ મળે છે : \VUgras{જંગલી
સસલું} \textsuperscript{(39)}, \VUgras{થ્રશ પક્ષી}
\textsuperscript{(40)}, \VUgras{તેતર} \textsuperscript{(41)},
\VUgras{બટેર} \textsuperscript{(42)}, \VUgras{જંગલી બતક}
\textsuperscript{(43)} કે \VUgras{ફિઝન્ટ} \textsuperscript{(44)} —
\VUgras{કેળાં} \textsuperscript{(46)} અને \VUgras{અનનાસ} જેવાં પરદેશી
ફળ જેટલી જ સહેલાઈથી.

%%K t15-09-1 p
9. --- કરિયાણાવાળાનો \VUgras{ગુમાસ્તો} \textsuperscript{(45)} ખૂબ
મદદરૂપ છે, અને જે માગો તે આપવા તૈયાર રહે છે : ગૂઝબેરીનો કે તેનો
\VUgras{મુરબ્બો} \textsuperscript{(47)}, \VUgras{ચરબી}
\textsuperscript{(48)}, \VUgras{નાની કાકડીઓ} \textsuperscript{(49)} કે
મીઠું ચડાવેલી \VUgras{સારડીન} \textsuperscript{(50)}.

%%K t15-09-2 p
એ \VUgras{ગ્રાહક બહેનો} \textsuperscript{(51)} માટે ખૂબ સગવડભર્યું છે,
જેમને રોજિંદા માલની બાજુમાં જ સૌથી દુર્લભ પહેલી ફસલ અને સૌથી સ્વાદિષ્ટ
ફળ મળે છે : રસદાર \VUgras{નાસપતી} \textsuperscript{(52)},
\VUgras{આલુબુખારા} \textsuperscript{(53)}, \VUgras{દ્રાક્ષ}
\textsuperscript{(54)}, \VUgras{આલૂ} \textsuperscript{(55)},
\VUgras{બદામ} \textsuperscript{(56)} અને \VUgras{અખરોટ}
\textsuperscript{(57)}.

%%K t15-tit-5 sub
\VUsaut{4.0mm}
\VUpk{10.2pt}{40}{ગ્રાહકો.}
\VUsaut{3.0mm}

%%K t15-10-1 p
10. --- \VUgras{ચોખા} \textsuperscript{(58)} અને \VUgras{મસૂર}
\textsuperscript{(59)} ધુમાડે પકવેલી \VUgras{હેરિંગની}
\textsuperscript{(60)} પેટી પાસે છે ; \VUgras{બટાટા}
\textsuperscript{(61)} અને \VUgras{વાલ} \textsuperscript{(63)}
\VUgras{સફરજન} \textsuperscript{(62)}, \VUgras{અંજીર}
\textsuperscript{(64)}, \VUgras{સૂકાં આલુબુખારા} \textsuperscript{(65)}
અને \VUgras{હેઝલનટની} \textsuperscript{(66)} બાજુમાં છે. એ દુકાનમાં
તાજાં \VUgras{ઈંડાં} \textsuperscript{(67)} અને \VUgras{ઑલિવ}
\textsuperscript{(68)} મળે છે ; એટલે \VUgras{ટોપલા}
\textsuperscript{(70)} અને \VUgras{સામાનની જાળીઓ}
\textsuperscript{(73)} પલકવારમાં ભરાઈ જાય છે.

%%K t15-11-1 p
11. --- એ ઉત્તમ પેઢીની \VUgras{કૉફીએ} \textsuperscript{(72)}
\VUgras{નોકરાણીઓ} \textsuperscript{(69)} અને \VUgras{રસોયણોમાં} વાજબી
નામના મેળવી છે, કેમ કે ઘરડો \VUgras{કરિયાણાવાળો} \textsuperscript{(71)}
પોતે જ તેને ખૂબ કાળજીથી શેકે છે.

%%K t15-tit-6 sub
\VUsaut{4.0mm}
\VUpk{10.2pt}{40}{જોવા જેવું.}
\VUsaut{3.0mm}

%%K t15-12-1 p
12. --- બધા બજારમાં ખરીદી કરવા જ નથી આવતા. મંડપ સુધી લઈ જતી ગલીની
રમણીય અવરજવરમાં જાતજાતના માણસો દેખાય છે. એક મળતાવડા \VUgras{કસાઈના
છોકરાની} \textsuperscript{(75)} પાસે, જે પોતાની આગળ \VUgras{હાથલારી}
\textsuperscript{(74)} ધકેલે છે, ત્યાં બે રજા પરના સિપાહી પોતાનો ચમકતો
ગણવેશ બતાવવામાં ખૂબ ગર્વ અનુભવે છે : વાદળી \VUgras{ડોલમન કોટ} અને
\VUgras{પીંછાવાળી ટોપી} પહેરેલો \VUgras{હુસાર} \textsuperscript{(76)},
અને ફૂલેલી બાંય તથા માથે \VUgras{ફેઝ ટોપીવાળો} \VUgras{ઝુઆવ
હવાલદાર}.

%%K t15-13-1 p
13. --- \VUgras{પાંઉવાળો} \textsuperscript{(80)}, જે \VUgras{પાંઉની
દુકાનના} \textsuperscript{(78)} ઉંબરે ઊભો છે, તેણે હમણાં જ પોતાના
\VUgras{પાંઉનો} \textsuperscript{(79)} લોટ ગૂંદ્યો છે અને ભઠ્ઠીમાંથી
\VUgras{ક્રોસાં} \textsuperscript{(81)}, \VUgras{નાના પાંઉ}
\textsuperscript{(82)} અને \VUgras{ગોળ પાંઉ} \textsuperscript{(83)}
કાઢ્યા છે, જે શોકેસમાં છે.

%%K t15-14-1 p
14. --- પાંઉની દુકાનની ઉપર એક કુશળ \VUgras{કપડાં સીવનારી બહેન}
\textsuperscript{(84)} રહે છે, જે કેટલીક \VUgras{ભરતકામ કરનારી બહેનોને}
\textsuperscript{(85)} કામે રાખે છે. તેની પડોશમાં એક \VUgras{ધોઈ-ઇસ્ત્રી
કરનારી બહેન} \textsuperscript{(86)} રહે છે, જે \VUgras{કપડાં}
\textsuperscript{(88)} ધોઈ-સૂકવ્યા પછી તેમને કાંજી ચડાવે છે અને
\VUgras{ઇસ્ત્રીથી} \textsuperscript{(87)} લીસાં કરે છે.

%%K t15-tit-7 sub
\VUsaut{4.0mm}
\VUpk{10.2pt}{40}{માંસ, માછલી અને શાક.}
\VUsaut{3.0mm}

%%K t15-15-1 p
15. --- ભોંયતળિયે આવેલી \VUgras{માંસની દુકાનમાં}
\textsuperscript{(89)} બધી જાતનું \VUgras{માંસ} વેચાય છે —
\VUgras{ઘેટાનું માંસ} \textsuperscript{(90)}, \VUgras{ગૌમાંસ}
\textsuperscript{(94)} કે \VUgras{વાછરડાનું માંસ}
\textsuperscript{(92)} — \VUgras{ઘેટાની જાંઘ, બીફસ્ટેક, શેકેલું માંસ}
અને \VUgras{કટલેટ} તરીકે. \VUgras{કસાઈનો કારીગર}
\textsuperscript{(93)} એ બધું માંસ કાપે છે અને \VUgras{મજ્જાવાળાં
હાડકાં} \textsuperscript{(96)} અલગ મૂકે છે. માંસની દુકાનના બારણે એક
\VUgras{છીપ વેચનારી બહેન} \textsuperscript{(97)} છે, જે \VUgras{છીપ}
\textsuperscript{(98)} વેચે છે. કોઈ ઇચ્છે તો તે એ ખોલી પણ આપે છે.

%%K t15-16-1 p
16. --- એ બધા વેપારીઓ પરવાનો કે જગ્યાનો કર ભરે છે ; એટલે
\VUgras{હવાલદાર} \textsuperscript{(100)} નિર્દયપણે બિચારી \VUgras{શાક
ફેરી કરનારી બહેનોની} \textsuperscript{(99)} પાછળ પડે છે, જે પોતાની
લારીઓમાં \VUgras{ગાજર, મૂળા, સલગમ, લીક} કે \VUgras{શતાવરીની ઝૂડીઓ,
લીલા વટાણા, પાલક, ચૂકો, ક્રેસ} અને \VUgras{પાર્સલી} ફેરવીને તેમની
હરીફાઈ કરે છે.

%%K t15-tit-8 sub
\VUsaut{4.0mm}
\VUpk{10.2pt}{40}{હવાલદારથી ચેતજો !}
\VUsaut{3.0mm}

%%K t15-17-1 p
17. --- જે છોકરો \VUgras{લસણ} \textsuperscript{(101)} અને
\VUgras{ડુંગળી} વેચે છે તે સારી પેઠે જાણે છે કે તેને પણ બજારની પાસે
પોતાનો માલ વેચવાની છૂટ નથી. તે ભાગે છે, અને દોડતાં દોડતાં
\VUgras{કાંટાળા લોબસ્ટર} અને \VUgras{લોબસ્ટર} \VUgras{વેચનારના}
\textsuperscript{(102)} \VUgras{કરચલા}, \VUgras{નદીના ઝીંગા} અને
\VUgras{ઝીંગાનો} ટોપલો ઢોળી નાખવાનું જોખમ વહોરે છે.

%%K t15-18-1 p
18. --- બધા દોડાદોડ કરે છે અને ભારે ઘોંઘાટ મચાવે છે ; પણ એથી ફેરી
કરતા \VUgras{કાચવાળાનો} \textsuperscript{(103)} અવાજ સાફ સંભળાતો
અટકતો નથી, કે \VUgras{કોલસાવાળાની} \textsuperscript{(104)} બૂમ પણ, જે
હમણાં જ \VUgras{કોલસાની ગૂણ} \textsuperscript{(105)} પહોંચાડવા થોડી વાર
માટે પોતાનું ગાડું છોડીને ગયો છે.

\VUsaut{6.0mm}
\VUfilet{22mm}
